% Краткое сообщение Екатерине: в работе обнаружены две ошибки. В задании 9 на первом листе общий множитель был сокращён до полного разложения числителя на множители. В задании 6 на втором листе при умножении неравенства на отрицательное число не был правильно изменён знак неравенства и знак правой части. Остальные читаемые задания выполнены верно и поэтому в ревью не включены.
\documentclass[12pt,a4paper]{article}
\usepackage[a4paper,left=20mm,right=20mm,top=18mm,bottom=20mm]{geometry}
\usepackage[T2A]{fontenc}
\usepackage[utf8]{inputenc}
\usepackage[russian]{babel}
\usepackage{amsmath,amssymb,mathtools}
\usepackage{array,booktabs,longtable}
\usepackage{xcolor}
\usepackage{hyperref}
\usepackage{tikz}

\hypersetup{
  colorlinks=true,
  linkcolor=blue!55!black,
  urlcolor=blue!55!black,
  pdftitle={Ревью домашней работы},
  pdfauthor={Лёвин Артём Александрович}
}

\setlength{\parindent}{0pt}
\setlength{\parskip}{0.55em}
\renewcommand{\arraystretch}{1.25}
\setlength{\tabcolsep}{3pt}

\newcommand{\errlabel}{\textcolor{red}{Ошибка:}}
\newcommand{\keyidea}{\textcolor{blue!60!black}{Ключевая идея:}}
\newcommand{\resultlabel}{\textcolor{teal!60!black}{Итог:}}

\begin{document}

\begin{titlepage}
\thispagestyle{empty}
\begin{center}
\vspace*{0.23\textheight}
{\Huge\bfseries Ревью домашней работы\par}
\vspace{2.2cm}
{\Large Автор: Лёвин Артём Александрович\par}
\vspace{0.8cm}
{\large 5 октября 2026 года\par}
\vfill
\begin{tikzpicture}[x=1cm,y=1cm]
\draw[blue!45!black,line width=0.8pt]
  (-3.8,0) .. controls (-2.0,0.55) and (2.0,-0.55) .. (3.8,0);
\end{tikzpicture}
\end{center}
\end{titlepage}

\newpage
\section*{Сводная таблица}
\small
\begin{longtable}{>{\raggedright\arraybackslash}p{1.6cm}
                  >{\raggedright\arraybackslash}p{2.8cm}
                  >{\raggedright\arraybackslash}p{4.4cm}
                  >{\raggedright\arraybackslash}p{3.2cm}
                  >{\raggedright\arraybackslash}p{3.2cm}}
\toprule
\textbf{№ задания} & \textbf{Тема} & \textbf{Суть ошибки} & \textbf{Ваш ответ} & \textbf{Правильный ответ} \\
\midrule
\endfirsthead
\toprule
\textbf{№ задания} & \textbf{Тема} & \textbf{Суть ошибки} & \textbf{Ваш ответ} & \textbf{Правильный ответ} \\
\midrule
\endhead
\hyperref[task:01-9]{9 (лист 01)} &
Сокращение алгебраических дробей &
Множитель $a+4b$ сокращён, пока он не является множителем всего числителя. &
Не указан. &
$\dfrac{1}{a+3}$ при исходных ограничениях. \\
\midrule
\hyperref[task:02-6]{6 (лист 02)} &
Неравенства и умножение на отрицательное число &
При умножении неравенства на $-0{,}5$ неверно изменены знак неравенства и знак правой части. &
Получена неверная оценка. &
Выражение меньше $-2$. \\
\bottomrule
\end{longtable}
\normalsize

\newpage
\thispagestyle{empty}
\vspace*{\fill}
\begin{center}
{\Huge\bfseries Разбор ошибок}
\end{center}
\vspace*{\fill}

\newpage
\phantomsection\label{task:01-9}
\section*{Задание 9 (лист 01)}

\subsection*{Текст задания}
Упростить выражение
\[
\frac{a^2+4ab-3a-12b}{a^2-16b^2}\cdot\frac{a-4b}{a^2-9}.
\]

\subsection*{Ваше решение}
В записи правильно начато разложение знаменателя разности квадратов и правильно вынесен множитель $a$ из первых двух слагаемых числителя:
\[
a^2-16b^2=(a-4b)(a+4b),
\]
\[
a^2+4ab-3a-12b=a(a+4b)-3a-12b.
\]
Также записано разложение
\[
a^2-9=(a-3)(a+3).
\]
После этого в решении выполняется сокращение множителя $a+4b$ до того, как весь числитель разложен на множители. Окончательный ответ не указан.

\subsection*{Ваш ответ}
Не указан.

\subsection*{Ошибка}
\errlabel{} недопустимое сокращение множителя внутри суммы.

\subsection*{Где именно возникает ошибка}
Последний корректный шаг --- представление числителя в виде
\[
a(a+4b)-3a-12b.
\]
Первый неверный шаг --- попытка сократить $a+4b$ со знаменателем, когда числитель ещё является суммой и разностью нескольких слагаемых.

\subsection*{Почему это неверно}
Сокращать в дроби можно только общий \emph{множитель} всего числителя и всего знаменателя. Например, из
\[
\frac{(a+4b)(a-3)}{(a+4b)(a-4b)}
\]
множитель $a+4b$ действительно можно сократить, если он не равен нулю.

Но в выражении
\[
a(a+4b)-3a-12b
\]
скобка $a+4b$ пока стоит множителем только у первого слагаемого $a(a+4b)$. Она ещё не записана как множитель \emph{всего} числителя. Поэтому сокращение на этом этапе меняет значение выражения.

Чтобы получить общий множитель, нужно сначала сгруппировать последние два слагаемых:
\[
-3a-12b=-3(a+4b).
\]
Тогда весь числитель примет вид
\[
a(a+4b)-3(a+4b),
\]
и только теперь общий множитель $a+4b$ можно вынести за скобки.

\subsection*{Ключевая идея}
\keyidea{} перед сокращением нужно убедиться, что сокращаемое выражение является множителем всего числителя и всего знаменателя, а не только отдельного слагаемого.

\subsection*{Исправление от последнего правильного шага}
Продолжим с корректной записи
\[
a(a+4b)-3a-12b.
\]
Сначала преобразуем последние два слагаемых:
\[
-3a-12b=-3(a+4b).
\]
Получаем
\[
a(a+4b)-3(a+4b).
\]
Теперь выносим общий множитель $a+4b$:
\[
a(a+4b)-3(a+4b)=(a-3)(a+4b).
\]
После этого сокращение становится законным.

\subsection*{Полное правильное решение}
Сначала учитываем ограничения исходного выражения. Знаменатели не должны обращаться в ноль:
\[
a^2-16b^2\ne 0,\qquad a^2-9\ne 0.
\]
Следовательно,
\[
a\ne 4b,\qquad a\ne -4b,\qquad a\ne 3,\qquad a\ne -3.
\]

Разложим числитель первой дроби группировкой:
\begin{align*}
a^2+4ab-3a-12b
&=a(a+4b)-3(a+4b)\\
&=(a-3)(a+4b).
\end{align*}

Разложим знаменатели по формуле разности квадратов:
\[
a^2-16b^2=(a-4b)(a+4b),
\]
\[
a^2-9=(a-3)(a+3).
\]

Подставляем разложения в исходное выражение:
\[
\frac{(a-3)(a+4b)}{(a-4b)(a+4b)}\cdot
\frac{a-4b}{(a-3)(a+3)}.
\]

При исходных ограничениях множители $a+4b$, $a-4b$ и $a-3$ не равны нулю, поэтому их можно сократить:
\[
\frac{1}{a+3}.
\]

\subsection*{Проверка}
Возьмём допустимые значения $a=5$ и $b=0$. Тогда исходное выражение равно
\[
\frac{25-15}{25}\cdot\frac{5}{25-9}
=\frac{10}{25}\cdot\frac{5}{16}
=\frac18.
\]
Упрощённое выражение даёт
\[
\frac{1}{5+3}=\frac18.
\]
Значения совпадают.

\subsection*{Итог}
\resultlabel{} правильное упрощение:
\[
\boxed{\frac{1}{a+3}},
\]
при ограничениях исходного выражения
\[
a\ne \pm 4b,\qquad a\ne \pm 3.
\]

\subsection*{Задача на закрепление}
Упростите выражение
\[
\frac{x^2+5xy-2x-10y}{x^2-25y^2}\cdot
\frac{x-5y}{x^2-4}.
\]

\newpage
\phantomsection\label{task:02-6}
\section*{Задание 6 (лист 02)}

\subsection*{Текст задания}
Известно, что $a>12$. Оценить выражение
\[
-0{,}5a+4.
\]

\subsection*{Ваше решение}
Вы правильно начали с условия $a>12$ и указали умножение обеих частей неравенства на отрицательное число $-0{,}5$. В следующем переходе у Вас получено неравенство, в котором левая часть равна $-0{,}5a$, а правая часть записана как положительное число $6$; направление неравенства также выбрано неверно. Затем к обеим частям прибавлено $4$.

\subsection*{Ваш ответ}
Получена неверная оценка выражения.

\subsection*{Ошибка}
\errlabel{} неверное преобразование неравенства при умножении на отрицательное число.

\subsection*{Где именно возникает ошибка}
Последний корректный шаг:
\[
a>12.
\]
Первый неверный шаг возникает при умножении обеих частей на $-0{,}5$.

Здесь одновременно нужно выполнить два действия:
\begin{itemize}
\item число $12$ умножить на $-0{,}5$, получив $-6$;
\item изменить направление знака неравенства, потому что умножение выполняется на отрицательное число.
\end{itemize}

\subsection*{Почему это неверно}
Для неравенств действует правило: если обе части строгого неравенства умножить или разделить на отрицательное число, знак неравенства меняется на противоположный.

Из
\[
a>12
\]
после умножения на $-0{,}5$ получается не знак «больше», а знак «меньше»:
\[
-0{,}5a<-0{,}5\cdot 12.
\]
Правая часть равна
\[
-0{,}5\cdot 12=-6,
\]
поэтому правильный промежуточный результат:
\[
-0{,}5a<-6.
\]

\subsection*{Ключевая идея}
\keyidea{} при умножении неравенства на отрицательное число нужно одновременно изменить направление знака неравенства и аккуратно сохранить знак произведения в каждой части.

\subsection*{Исправление от последнего правильного шага}
Начинаем с
\[
a>12.
\]
Умножаем обе части на $-0{,}5$. Поскольку множитель отрицательный, знак меняется:
\[
-0{,}5a<-6.
\]
Теперь прибавляем $4$ к обеим частям. При сложении с одним и тем же числом направление неравенства не меняется:
\[
-0{,}5a+4<-6+4.
\]
Следовательно,
\[
-0{,}5a+4<-2.
\]

\subsection*{Полное правильное решение}
Дано:
\[
a>12.
\]
Умножаем обе части на $-0{,}5$ и меняем знак неравенства:
\[
-0{,}5a<-6.
\]
Прибавляем $4$ к обеим частям:
\[
-0{,}5a+4<-2.
\]
Значит, значение выражения $-0{,}5a+4$ всегда меньше $-2$.

\subsection*{Проверка}
Возьмём любое число, удовлетворяющее условию, например $a=14$. Тогда
\[
-0{,}5\cdot 14+4=-7+4=-3.
\]
Действительно,
\[
-3<-2.
\]
Проверка согласуется с полученной оценкой.

\subsection*{Итог}
\resultlabel{} при $a>12$ выполняется
\[
\boxed{-0{,}5a+4<-2}.
\]

\subsection*{Задача на закрепление}
Известно, что $c>18$. Оцените выражение
\[
-\frac{2}{3}c+5.
\]

\vfill
\begin{center}
\small Лёвин Артём Александрович эксклюзивно для Екатерины Скелли
\end{center}

\end{document}
